% Generated by scripts/materialize_unified_paper.py; do not edit.
\section{Query complexity and the classical landscape}\label{qua:sec:complexity}

For the size-biased recursion, a capped sample from a nonempty subset of the
$M$-element universe costs $O(\sqrt M)$ membership tests and
$O(\log(2r))$ repetitions give failure probability $O(1/r)$.  The predicate at
level $i$ costs $O(i+1)$ edge queries, hence one fixed-cap run costs
\[
 O\!\left(\sqrt M\log(2r)\sum_{i=0}^{r-1}(i+1)\right)
 =O(K^2 2^K\log K).
\]
The repetition in the preceding section proves the following preliminary
bound in worst-case, rather than merely expected, query complexity.

\begin{proposition}[Size-biased quantum recursion]
\label{qua:prop:size-biased-algorithm}
Under the hypotheses of \cref{qua:thm:main}, the estimation-free algorithm returns
a verified homogeneous $K$-set with probability at least $1-\eta$ using
\[
 O\!\left(K^2 2^K\log K\log\frac1\eta\right)
\]
edge queries.
\end{proposition}

We next analyze the sharper scale-aware recursion.

By \cref{qua:lem:survival}, a capped search block at level $i$ uses
$O(2^{i/2})$ membership tests.  Computing and uncomputing membership and
checking the sampled edge costs $O(i+1)$ edge queries.  Put
\[
 b_i=(i+1)2^{i/2}.
\]
The batch cap uses $O(m_i+\log(r/\eta))=O(m_i)$ blocks.  Hence its total
sampling cost is
\begin{equation}\label{qua:eq:query-sum}
 O\!\left(
 \log\frac r\eta
 \sum_{i=0}^{r-1}b_i\varepsilon_i^{-2}
 \right).
\end{equation}

There are at most six indices at each value of $d_i$, so
\[
 Z=\sum_i2^{-d_i}<6\sum_{d\ge0}2^{-d}=12.
\]
Using \cref{qua:eq:error-schedule},
\[
 \sum_{i=0}^{r-1}b_i\varepsilon_i^{-2}
 =O\!\left(
 \sum_{i=0}^{r-1}(i+1)2^{i/2}4^{d_i}
 \right).
\]
Write $h=r-1-i$.  Since
$4^{\lceil h/6\rceil}\le4\cdot2^{h/3}$, the last display is bounded by
\[
 O\!\left(
 2^{r/2}\sum_{h=0}^{r-1}(r-h)2^{-h/6}
 \right)
 =O(r2^{r/2}).
\]
The high-confidence pivot searches, final search, and $K^2$ verification
fit the same bound.  Since $r=2K-3$, this proves \cref{qua:thm:main}.

The gate complexity is
\[
 2^K\operatorname{poly}\!\left(
 K,\log N,\log\frac1\eta
 \right)
\]
plus the gate cost of one supplied edge-oracle evaluation.  No QRAM is
assumed.  The $O(K)$ measured pivot labels and colours are classical controls
for reversible comparisons, conjunctions, and uncomputation.

The known classical lower bound remains below the quantum upper bound.  The
following direct consequence of the random-Painter analysis of
Conlon--Fox--Grinshpun--He \cite{ConlonFoxGrinshpunHe2019} records the best
bound we obtain from the current literature; details appear in
\cref{qua:app:classical-lower}.

\begin{proposition}[Randomized classical lower bound]
\label{qua:prop:classical-lower}
For $K\ge8$ and $M=4^{K-1}$, every randomized edge-query algorithm that finds
a homogeneous $K$-set with worst-case probability at least $2/3$ uses
\[
 \Omega\!\left(2^{(2-\sqrt2)K}\right)
 =\Omega\!\left(M^{1-1/\sqrt2}\right)
\]
queries.
\end{proposition}

The exponent $1-1/\sqrt2\approx0.292893$ does not reach the quantum upper
exponent $1/2$, but it strengthens the earlier $M^{1/4}$ black-box lower
bound recorded by Komargodski--Naor--Yogev in this parameterization
\cite{KomargodskiNaorYogev2017}.  Therefore no
quantum-versus-best-randomized-classical
separation is claimed.  A randomized lower bound
$2^{(1+\delta)K}$ for any fixed $\delta>0$ would, through the online Ramsey
connection, improve the long-standing exponential lower bound for diagonal
Ramsey numbers.  The missing separation is tied to a major combinatorial
barrier rather than a routine query lemma.

The size-biased mechanism also extends to any fixed number of edge colours.
The statement and its one-line recurrence proof are given in
\cref{qua:app:multicolour}; this extension is not needed for the sharper binary
theorem.
